%% 5-14 ふしぎな数 e（第11回：ページ制限なし。X17：複利の「半年で50％」を、1年あたりの割合から1段ずつ説明する）
\bkchapter{5}{14}{ふしぎな数 $e$}{\bklevelrange{3}{4}}{ふしぎな数 e}
\bkmichishirube{3-11,3-12,5-6}{5-17,6-4,9-5}

\begin{goal}
\begin{itemize}[label={\color{bkTheme}●}, leftmargin=1.6\zw, itemsep=0pt, topsep=1pt]
\item \bklevel{3}\enspace $e$ が $\Bigl(1+\dfrac{1}{n}\Bigr)^n$ の近づく先であることを、表計算で確かめられる。
\item \bklevel{3}\enspace $(e^{ax})'=ae^{ax}$ を使って、ふえる量・へる量の変化の速さを求められる。
\item \bklevel{4}\enspace 自然対数 $\ln x$ を微分・積分できる。\bkidx{たいすう}{対数}
\end{itemize}
\end{goal}

\begin{hatena}
\noindent\begin{minipage}[c]{0.68\linewidth}\raggedright
利率が「1年あたり100\,\%」という、ゆめのような銀行があります。1年たつと、預けたお金と同じだけの利息（ふえる分）がつきます。1万円を預けると、1年後に2万円です。

この銀行が、利息を半年ごとに計算してくれることになりました。1年あたりの利率は100\,\%のままです。1年後のお金は、2万円より多くなるでしょうか。毎月、毎日……と、もっと細かく計算してもらえば、いくらでも大きくなるのでしょうか。
\end{minipage}\hfill
\begin{minipage}[c]{0.3\linewidth}\centering
\begin{tikzpicture}[x=1mm, y=1mm]
  \draw[bkInk, line width=0.6pt] (-1,0) -- (36,0);
  \fill[bkYellow!60] (0,0) rectangle (8,12);
  \draw[bkYellow!70!black, line width=0.6pt] (0,0) rectangle (8,12);
  \node[font=\sffamily\scriptsize, anchor=south] at (4,13) {1万円};
  \foreach \px in {13,26} {\draw[bkGray, line width=0.6pt, dash pattern=on 2pt off 1.5pt] (\px,0) rectangle (\px+8,22); \node[font=\sffamily\bfseries\small, text=bkGray] at (\px+4,11) {?};}
  \node[font=\sffamily\scriptsize, anchor=north] at (4,-1) {はじめ};
  \node[font=\sffamily\scriptsize, anchor=north] at (17,-1) {半年};
  \node[font=\sffamily\scriptsize, anchor=north] at (30,-1) {1年};
\end{tikzpicture}
\end{minipage}
\end{hatena}

この章では、半年ごとの計算を1段ずつ追いかけ、その先にあらわれる数 $e$（ネイピア数）を調べます。$e$ は、放射線のへり方、お湯の冷め方、電気回路など、自然のあらゆる所に顔を出します。

%% ---------------- やってみよう・思い出そう ----------------
\begin{yattemiyou}[半年ごとに利息をつけてもらおう]
「1年あたり100\,\%」の利率で、1万円を預けます。銀行は、半年ごとに利息を計算して、預けているお金にたしてくれます。電卓を使ってよいです。
\begin{enumerate}[label={\bkcircled[bkGreen]{\arabic*}}, leftmargin=2.4\zw, itemsep=1pt, topsep=3pt]
\item 半年は、1年の何分の1かな？　半年分の利率は、100\,\%の何分の1になるかな？
\item はじめの半年：1万円に、半年分の利息をつけよう。半年後にはいくら？
\item 次の半年：半年後のお金に、また半年分の利息をつけよう。1年後にはいくら？
\item 1年に1回だけ計算したとき（2万円）とくらべよう。差は、どこから出たのかな？
\end{enumerate}
\begin{bkcenter}
\begin{bktable}{colspec={Q[l,m,wd=40mm] *{3}{Q[c,m,wd=20mm]}}, rowsep=1.5pt}
 & はじめ & 半年後 & 1年後 \\
1年に1回計算する（万円） & 1 & 1 & 2 \\
半年ごとに計算する（万円） & 1 & ？ & ？ \\
\end{bktable}
\end{bkcenter}
\end{yattemiyou}

\noindent ①半年は1年の $\dfrac{1}{2}$ なので、半年分の利率は $100\,\%\times\dfrac{1}{2}=50\,\%$ です。②1万円の50\,\%は0.5万円なので、半年後は $1+0.5=1.5$ 万円。③次の半年も50\,\%です。こんどは1.5万円に50\,\%がつくので、利息は $1.5\times0.5=0.75$ 万円、1年後は $1.5+0.75=2.25$ 万円です。④1年に1回のときより、0.25万円多くなりました。

\begin{bkfigure}
\begin{tikzpicture}[x=1mm, y=1mm]
  \def\u{16}% 1万円の高さ（mm）
  \draw[bkInk, line width=0.7pt] (-4,0) -- (98,0);
  % はじめ
  \fill[bkYellow!60] (0,0) rectangle (14,\u); \draw[bkYellow!70!black, line width=0.6pt] (0,0) rectangle (14,\u);
  % 半年
  \fill[bkYellow!60] (40,0) rectangle (54,\u); \fill[bkOrange!55] (40,\u) rectangle (54,1.5*\u);
  \draw[bkYellow!70!black, line width=0.6pt] (40,0) rectangle (54,\u); \draw[bkOrange!80!black, line width=0.6pt] (40,\u) rectangle (54,1.5*\u);
  % 1年
  \fill[bkYellow!60] (80,0) rectangle (94,\u); \fill[bkOrange!55] (80,\u) rectangle (94,1.5*\u);
  \fill[bkBlue!35] (80,1.5*\u) rectangle (94,2*\u); \fill[bkRed!55] (80,2*\u) rectangle (94,2.25*\u);
  \draw[bkYellow!70!black, line width=0.6pt] (80,0) rectangle (94,\u); \draw[bkOrange!80!black, line width=0.6pt] (80,\u) rectangle (94,1.5*\u);
  \draw[bkBlue!80!black, line width=0.6pt] (80,1.5*\u) rectangle (94,2*\u); \draw[bkRed!80!black, line width=0.6pt] (80,2*\u) rectangle (94,2.25*\u);
  % 合計
  \node[font=\sffamily\small, anchor=south] at (7,\u+1.5) {1万円};
  \node[font=\sffamily\small, anchor=south] at (47,1.5*\u+1.5) {1.5万円};
  \node[font=\sffamily\bfseries\small, anchor=south] at (87,2.25*\u+1.5) {2.25万円};
  % 50％ふえる（×1.5）
  \foreach \a/\b in {16/38, 56/78} {
    \node[font=\sffamily\scriptsize, align=center, anchor=south, text=bkInk] at ({(\a+\b)/2},8.5) {50\,\%ふえる\\（×1.5）};
    \draw[bk arrow, bkGray, line width=0.8pt] (\a,6) -- (\b,6);}
  % 時間の目もり
  \node[font=\sffamily\small, anchor=north] at (7,-1.5) {はじめ};
  \node[font=\sffamily\small, anchor=north] at (47,-1.5) {半年};
  \node[font=\sffamily\small, anchor=north] at (87,-1.5) {1年};
  % 凡例
  \begin{scope}[shift={(0,-14)}]
    \foreach \c/\t/\px/\py in {bkYellow!60/はじめの1万円/0/0, bkOrange!55/1回めの利息0.5万円/52/0, bkBlue!35/2回めの利息のうち、1万円の分0.5万円/0/-6, bkRed!55/2回めの利息のうち、利息の分0.25万円/0/-12} {
      \fill[\c] (\px,\py-1.6) rectangle ++(3.2,3.2);
      \node[font=\sffamily\scriptsize, anchor=west, inner sep=0pt] at (\px+4.8,\py) {\t};}
  \end{scope}
\end{tikzpicture}
\bkfigcaption{半年ごとに計算すると、2回めの利息は、1万円だけでなく、1回めの利息0.5万円にもつく（赤）。}
\end{bkfigure}

\noindent 1年に1回のときとの差の0.25万円（赤）は、1回めの利息0.5万円に、2回めの半年でついた50\,\%の利息（$0.5\times0.5=0.25$）です。利息にも利息がつくのです。このような利息のつけ方を\textbf{複利}（ふくり）といいます。\bkidx{ふくり}{複利}

\begin{omoidasou}
\bkrecall{2-5}{50\,\%ふえる＝もとの1倍＋0.5倍＝1.5倍。20\,\%引き＝0.8倍と同じ考え。}
\bkrecall{2-7}{同じ数を何回もかけるのが累乗。\\ $1.5\times1.5=1.5^2$、\quad $1.25\times1.25\times1.25\times1.25=1.25^4$}
\bkrecall{3-11}{指数関数 $y=a^x$。$a^{x+h}=a^x\times a^h$、$(a^m)^n=a^{mn}$、$a^0=1$。たとえば $2^{3+1}=2^3\times2^1=8\times2=16$}
\bkrecall{3-12}{$\log_a b$ は「$a$ を何乗すると $b$ になるか」で、$a^{\log_a b}=b$。\\ $\log_a(bc)=\log_a b+\log_a c$、\quad $\log_a b^r=r\log_a b$}
\bkrecall{5-6}{合成関数の微分：外側を微分して、内側の微分をかける。}
\bkrecall{5-8}{$h$ が小さいとき、$f(a+h)\approx f(a)+f'(a)h$（接線で近似）}
\end{omoidasou}

\clearpage
%% ---------------- 数学の道具 ----------------
\bksection{doguu}{数学の道具}

\bksubsection{1}{「1年あたり100\,\%」を、半年に分ける}

\label{bk:5-14:bunkatsu}「1年あたり100\,\%」の「あたり」は、時速の「あたり」と同じです。時速\qty{60}{km}は「1時間あたり\qty{60}{km}」でした。30分（0.5時間）走ると、
\[ \qty{60}{km/h}\times\qty{0.5}{h}=\qty{30}{km} \]
進みます。時間が半分なら、進む距離も半分です。利率も同じで、「1年あたり100\,\%」の利率で半年（0.5年）預けると、つく利息は
\[ 100\,\%\times0.5=50\,\% \]
です。\textbf{時間が半分なので、利息の割合も半分の50\,\%}になります。

\noindent\textbf{なぜ「1.5倍」なのか}\quad 50\,\%の利息がつくと、お金は「もとの1倍」に「0.5倍の利息」をたした $1+0.5=1.5$ 倍になります（2-5）。だから、半年ごとに1.5倍が2回で
\[ 1\ \text{万円}\times1.5\times1.5=1\ \text{万円}\times1.5^2=2.25\ \text{万円} \]
です。

\noindent\textbf{まちがえやすい考え方1}\quad 「はじめの半年は50\,\%、次の半年は100\,\%」と考えると、$1.5\times2=3$ 万円になってしまいます。100\,\%は「1年あたり」の割合でした。2回めも半年（0.5年）なので、2回めの利息も50\,\%です。半年ごとの計算では、毎回50\,\%です。

\noindent\textbf{まちがえやすい考え方2}\quad 「50\,\%と50\,\%で、合わせて100\,\%なので2万円」とすると、2回めの利息が1万円だけにつくと考えたことになります。でも、2回めの半年が始まるときには、預けているお金はもう1.5万円です。2回めの利息は、1.5万円の50\,\%（0.75万円）です。

\bksubsection{2}{もっと細かく分けると：$\Bigl(1+\dfrac{1}{n}\Bigr)^n$}

\label{bk:5-14:n}\textbf{3か月ごと（1年に4回）}\quad 3か月は1年の $\dfrac{1}{4}$ なので、1回分の利率は $100\,\%\times\dfrac{1}{4}=25\,\%$、1回ごとに $1+0.25=1.25$ 倍です。

\begin{bkcenter}
\begin{bktable}{colspec={Q[l,m,wd=22mm] Q[c,m,wd=38mm] Q[c,m,wd=30mm]}, rowsep=1.5pt, row{1}={bg=bkTheme!10}}
時期 & 計算 & お金（万円） \\
はじめ & & 1 \\
3か月後 & $1\times1.25$ & 1.25 \\
6か月後 & $1.25\times1.25$ & 1.5625 \\
9か月後 & $1.5625\times1.25$ & 1.953125 \\
1年後 & $1.953125\times1.25$ & 2.44140625 \\
\end{bktable}
\bktabcaption{3か月ごとに1.25倍。1年後は $1.25^4=2.4414\cdots$ 倍。}
\end{bkcenter}

\noindent\textbf{1年を $n$ 回に分ける}\quad 同じように考えると、
\begin{enumerate}[label={\bkcircled[bkTheme]{\arabic*}}, leftmargin=2.4\zw, itemsep=1pt, topsep=2pt]
\item 1回分の時間は1年の $\dfrac{1}{n}$ なので、1回分の利率は $100\,\%\times\dfrac{1}{n}$、つまり $\dfrac{1}{n}$ です（100\,\%＝1）。
\item 1回ごとに、お金は $\Bigl(1+\dfrac{1}{n}\Bigr)$ 倍になります。
\item それを $n$ 回くり返すので、1年後には $\Bigl(1+\dfrac{1}{n}\Bigr)^n$ 倍です。指数の $n$ は「$n$ 回かける」の $n$ です（2-7）。
\end{enumerate}
表計算で、$n$ を大きくしていきます。
\begin{bkcenter}
\bksheet{3}{28}{%
1 & ふやし方 & $n$ & $\left(1+\frac{1}{n}\right)^n$ \\
2 & 1年に1回 & 1 & 2 \\
3 & 半年ごと & 2 & 2.25 \\
4 & 3か月ごと & 4 & 2.4414… \\
5 & 毎月 & 12 & 2.6130… \\
6 & 毎日 & 365 & 2.7145… \\
7 & （細かく） & 10000 & 2.71814… \\
8 & （もっと細かく） & 1000000 & 2.718280… \\}
\bktabcaption{$n$ を大きくすると、2.71828… に近づいていく。}
\end{bkcenter}

\begin{bkfigure}
\begin{tikzpicture}
\begin{axis}[bkplot, width=100mm, height=46mm, xmin=0, xmax=52, ymin=1.8, ymax=3, xtick={1,10,20,30,40,50}, ytick={2,2.25,2.5,2.718}, yticklabels={2,2.25,2.5,$e$}, xlabel={分ける回数 $n$}, ylabel={1年後（万円）}]
  \addplot[bkthin, dashed, domain=0:52] {2.71828};
  \addplot[only marks, mark=*, mark size=1.3pt, bkBlue, samples at={1,2,...,50}] {(1+1/x)^x};
\end{axis}
\end{tikzpicture}
\bkfigcaption{$n$ をふやすと $\bigl(1+\frac{1}{n}\bigr)^n$ はふえていくが、点線の $e=2.718\cdots$ をこえない。}
\end{bkfigure}

細かく分けるほどお金はふえますが、ふえ方はだんだん小さくなります。1回から2回で0.25、2回から4回で約0.19、4回から12回で約0.17ふえますが、365回から1000000回にしても0.004ほどしかふえません。1回分の利息がとても小さくなるので、「利息の利息」も、ほんの少ししかふえないのです。$n$ をどんなに大きくしても3にもとどかず、2.71828… に近づいていきます。この近づく先を $e$ と書きます。
\[ e=\lim_{n\to\infty}\Bigl(1+\frac{1}{n}\Bigr)^n=2.718281828\cdots \]
$e$ は、$\pi$ と同じように、小数がどこまでも続き、くり返さない数です。これが\bkhatenaref の答えです。いくら細かく計算してもらっても、1年後のお金は約2.718万円より大きくはなりません。\bkidx{いー}{$e$（ネイピア数）}

\noindent\textbf{なぜ細かく分けて考えるの？}\quad お金のふえ方は銀行の決まりしだいですが、自然の変化は「1年に1回」まとめて起きるわけではありません。細菌は、いまいる数に比例して、いつも少しずつふえています。お湯は、いつも少しずつ冷めています。分け方をどんどん細かくしていくと、「いつも少しずつ、いまある量に比例してふえる」ようすに近づきます。その行き着く先にあらわれるのが $e$ です。このことを、③からの微分で確かめます。

\bksubsection{3}{指数関数を微分すると、自分に比例する}

\label{bk:5-14:shisuu}$f(x)=a^x$ を微分してみましょう。5-3の決め方どおり、$\dfrac{f(x+h)-f(x)}{h}$ を計算して、$h$ を0に近づけます。指数のきまり（3-11）
\[ a^{x+h}=a^x\times a^h\qquad（たとえば\ 2^{3+1}=2^3\times2^1=8\times2=16） \]
を使うと、
\[ \frac{a^{x+h}-a^x}{h}=\frac{a^x\times a^h-a^x}{h}=a^x\times\frac{a^h-1}{h} \]
最後の式は、分子の2つの項に共通な $a^x$ をくくり出しました。右の $\dfrac{a^h-1}{h}$ には $x$ が入っていません。だから、$a$ が決まれば、どの $x$ でも同じ数です。$h$ を0に近づけたときの近づく先を $k_a$ と書くと、
\[ (a^x)'=k_a\,a^x \]
\textbf{指数関数を微分すると、自分自身の $k_a$ 倍になります}。「いまある量に比例してふえる」（3-11）ことの、微分による表し方です。また、$x=0$ を入れると $a^0=1$ なので $f'(0)=k_a$。$k_a$ は $x=0$ での接線の傾きです。

\noindent\textbf{数で確かめる}\quad $a=2$ で、$h=0.0001$ とします。電卓で $2^{0.0001}=1.0000693\cdots$ なので、
\[ \frac{2^{0.0001}-1}{0.0001}=\frac{0.0000693\cdots}{0.0001}=0.693\cdots \]
$y=2^x$ の各点の傾きと高さをくらべると、どこでも「傾き÷高さ」が0.693です。
\begin{bkcenter}
\begin{bktable}{colspec={Q[l,m,wd=34mm] *{4}{Q[c,m,wd=14mm]}}, rowsep=1.5pt}
$x$ & 0 & 1 & 2 & 3 \\
高さ $2^x$ & 1 & 2 & 4 & 8 \\
傾き & 0.693 & 1.386 & 2.773 & 5.545 \\
傾き÷高さ & 0.693 & 0.693 & 0.693 & 0.693 \\
\end{bktable}
\bktabcaption{$y=2^x$ の傾きは、いつも高さの0.693倍。}
\end{bkcenter}

ほかの底についても、$k_a$ を表計算で $h=0.0001$ として調べました。
\begin{bkcenter}
\begin{bktable}{colspec={Q[c,m,wd=24mm] *{5}{Q[c,m,wd=14mm]}}, rowsep=1.5pt}
底 $a$ & 2 & 2.5 & 2.7 & 2.8 & 3 \\
$k_a\approx\dfrac{a^h-1}{h}$ & 0.693 & 0.916 & 0.993 & 1.030 & 1.099 \\
\end{bktable}
\bktabcaption{底が2.7と2.8のあいだで、$k_a$ がちょうど1になりそう。}
\end{bkcenter}

\bksubsection{4}{$e$ は「微分しても変わらない」底}

\label{bk:5-14:e}$k_a=1$ になる底を見つければ、その指数関数は、微分しても自分自身のままです。その底が、ちょうど②の $e=2.71828\cdots$ です。$a=e$ で計算すると、$\dfrac{e^{0.0001}-1}{0.0001}=1.00005$ で、たしかに1に近くなります。
\[ (e^x)'=e^x \]

\noindent\textbf{なぜ利息の $e$ と同じなのか}\quad $k_e=1$ は、$y=e^x$ の $x=0$ での接線が、点 $(0,\,1)$ を通る傾き1の直線 $y=1+x$ だということです。接点の近くでは、グラフと接線はほとんど重なるので、小さい $h$ では（5-8）
\[ e^h\approx1+h \]
たとえば $e^{0.1}=1.10517\cdots$ と $1.1$、$e^{0.01}=1.01005\cdots$ と $1.01$ です。ここで、$n$ を大きな数として、$h$ に $\dfrac{1}{n}$ を入れると、
\[ e^{\frac{1}{n}}\approx1+\frac{1}{n} \]
です。両辺を $n$ 乗します。左辺は、$(a^m)^n=a^{mn}$（3-11）を使うと
\[ \Bigl(e^{\frac{1}{n}}\Bigr)^n=e^{\frac{1}{n}\times n}=e^1=e \]
なので、
\[ e\approx\Bigl(1+\frac{1}{n}\Bigr)^n \]
$n$ が大きいほど近似は正確になるので、②の表の数が $e$ に近づいたのです。

\begin{bkfigure}
\begin{tikzpicture}
\begin{axis}[bkplot, width=86mm, height=52mm, xmin=-2.2, xmax=2.3, ymin=0, ymax=7.6, xtick={-2,-1,1,2}, ytick={2,4,6}, xlabel={$x$}, ylabel={$y$}, yticklabel style={fill=white, inner sep=1pt, outer sep=2.3pt}]
  \addplot[bkthin, domain=-2.1:2.1] {2^x};
  \addplot[bkthin, dashed, domain=-2.1:1.75] {3^x};
  \addplot[bkcurve, domain=-2.1:2] {exp(x)};
  \addplot[bkcurve2, line width=1pt, domain=-1:1.6] {1+x};
  \node[font=\sffamily\scriptsize, anchor=west] at (axis cs:2.2,4.29) {$2^x$};
  \node[font=\sffamily\scriptsize, anchor=east] at (axis cs:1.6,6.9) {$3^x$};
  \node[font=\sffamily\scriptsize, anchor=west, text=bkBlue] at (axis cs:2.08,7.3) {$e^x$};
  % 格子の線が「＋」の上を通らないように、空いている左上（x<0、y=4〜6）に置く
  \node[font=\sffamily\scriptsize, align=center, text=bkOrange!80!black, fill=white, inner sep=1pt] at (axis cs:-1.5,5) {$y=1+x$\\（傾き1）};
\end{axis}
\end{tikzpicture}
\bkfigcaption{$y=e^x$ は、$(0,\,1)$ での接線の傾きがちょうど1になる指数関数。$2^x$ はそれよりゆるく、$3^x$ は急。}
\end{bkfigure}

\begin{seikaku}
\textbf{$e$ の決め方}\par
この本では、$e$ を「$k_a=1$ になる底」、つまり $\displaystyle\lim_{h\to0}\frac{e^h-1}{h}=1$ となる数として決めました。大学の教科書では、$e=\displaystyle\lim_{n\to\infty}\Bigl(1+\frac{1}{n}\Bigr)^n$ で $e$ を決め、そこから $(e^x)'=e^x$ を導くことが多いです。どちらから始めても、同じ数 $2.71828\cdots$ になり、同じ性質がすべて導けます。また、$e$ は分数で表せない数（2-9）であることが証明されています。
\end{seikaku}

\bksubsection{5}{$e^{ax}$ の微分}

\label{bk:5-14:eax}合成関数の微分（5-6）を使います。内側を $u=ax$、外側を $e^u$ とします。外側 $e^u$ を $u$ で微分すると $e^u$（④）、内側 $ax$ を $x$ で微分すると $a$ です。2つをかけて、
\[ (e^{ax})'=e^{ax}\times a=a\,e^{ax} \]
たとえば $a=2$ なら $(e^{2x})'=2e^{2x}$ で、$x=0$ での傾きは2です。表計算で $\dfrac{e^{0.0002}-1}{0.0001}=2.0002$ と確かめられます。

$a>0$ なら、$e^{ax}$ はふえていき、ふえる速さはいつも自分の $a$ 倍です。$a<0$ ならへっていき、へる速さが自分に比例します。放射線を出す物質の量 $N$ が $N=N_0e^{-\lambda t}$（$\lambda$ はラムダ）でへるのは、「へる速さが、いまある量に比例する」からです。\bkidx{しすうかんすう}{指数関数の微分}

\noindent\textbf{例}\quad $(e^{3x})'=3e^{3x}$、\quad $(5e^{-2t})'=5\times(-2)e^{-2t}=-10e^{-2t}$\par
\noindent\hspace*{2.5\zw}$(e^{x^2})'=e^{x^2}\times(x^2)'=e^{x^2}\times2x$（内側を $u=x^2$ とする）

\bksubsection{6}{自然対数 $\ln$}

\label{bk:5-14:ln}$e$ を底とする対数 $\log_e x$ を\textbf{自然対数}といい、$\ln x$ と書きます（エル・エヌ・エックス。natural logarithm の略）。「$e$ を何乗すると $x$ になるか」です（3-12）。\bkidx{しぜんたいすう}{自然対数 $\ln$}
\[ e^{\ln x}=x,\qquad \ln e^{y}=y,\qquad \ln1=0,\qquad \ln e=1 \]
$\ln1=0$ なのは、$e^0=1$（0乗すると1）だからです。$\ln e=1$ なのは、$e^1=e$ だからです。

たとえば $\ln2=0.693\cdots$ は、「$e$ を約0.693乗すると2になる」ということです。電卓で確かめると、
\[ 2.71828^{0.693}=1.9997\cdots\approx2 \]
です。電卓や表計算では、\texttt{LN(2)} で求められます。

\noindent\textbf{どんな指数関数も $e$ で書ける}\par
$a=e^{\ln a}$ なので、3-11の $(a^m)^n=a^{mn}$ を使うと
\[ a^x=(e^{\ln a})^x=e^{x\ln a} \]
です。たとえば $2=e^{0.693}$ なので、$2^x=e^{0.693x}$ です。これを、⑤の $(e^{ax})'=ae^{ax}$ の $a$ のかわりに $\ln a$ として微分すると、
\[ (a^x)'=(e^{x\ln a})'=\ln a\times e^{x\ln a}=a^x\ln a \]
③の表の $k_a$ は、$\ln a$ だったのです。たしかに
\[ k_2=0.693=\ln2,\qquad k_3=1.099=\ln3 \]
です。

\bksubsection{7}{$\ln x$ の微分：グラフを折り返して考える}

\label{bk:5-14:lnbibun}$y=\ln x$ は、「$x=e^y$」と同じことです。$x$ と $y$ の役割を入れかえた関係なので、$y=\ln x$ のグラフは、$y=e^x$ のグラフを、直線 $y=x$ で折り返したものです（3-12）。

\noindent\begin{minipage}[c]{0.41\linewidth}
\centering
\begin{tikzpicture}
\begin{axis}[bkplot, width=56mm, height=56mm, xmin=-1.3, xmax=4.4, ymin=-1.3, ymax=4.4, xtick={1,2,3,4}, ytick={1,2,3,4}, xticklabels={,2,3,4}, yticklabels={,2,3,4}, xlabel={$x$}, ylabel={$y$}, tick label style={fill=white, inner sep=1pt, outer sep=2.3pt}]
  \addplot[bkthin, dashed, domain=-1.3:4.2] {x};
  \addplot[bkcurve2, domain=-1.3:1.45] {exp(x)};
  \addplot[bkcurve, domain=0.28:4.3] {ln(x)};
  \addplot[bkcurve2, line width=0.8pt, domain=0.15:1.25] {2+2*(x-0.693)};
  \addplot[bkcurve, line width=0.8pt, domain=1:3.2] {0.693+0.5*(x-2)};
  \addplot[only marks, mark=*, mark size=1.4pt, bkInk] coordinates {(0.693,2) (2,0.693)};
  \node[font=\sffamily\scriptsize, text=bkOrange!80!black, anchor=west, fill=white, inner sep=1pt] at (axis cs:1.6,4.0) {$e^x$};
  \node[font=\sffamily\scriptsize, text=bkBlue, anchor=south, fill=white, inner sep=1pt] at (axis cs:3.9,1.75) {$\ln x$};
\end{axis}
\end{tikzpicture}\par\vspace{1pt}
{\sffamily\scriptsize\color{bkOrange!80!black}点 $(\ln2,\,2)$ での傾き2}\par
{\sffamily\scriptsize\color{bkBlue}点 $(2,\,\ln2)$ での傾き $\frac{1}{2}$}
\end{minipage}\hfill
\begin{minipage}[c]{0.56\linewidth}
折り返すと、横の長さとたての長さが入れかわります。だから、接線の傾き（たて÷横）は、逆数になります。$y=e^x$ の点 $(\ln2,\,2)$ での傾きは $e^{\ln2}=2$ なので、折り返した $y=\ln x$ の点 $(2,\,\ln2)$ での傾きは $\dfrac{1}{2}$ です。

どの点でも同じです。$y=e^x$ の点 $(p,\ e^p)$ での傾きは $e^p$。折り返すと、この点は $y=\ln x$ の点 $(e^p,\ p)$ になり、傾きは $\dfrac{1}{e^p}$。この点の $x$ 座標は $e^p$ なので、傾きは $\dfrac{1}{x}$ です。
\end{minipage}
\[ (\ln x)'=\frac{1}{x}\qquad（x>0） \]
\noindent\textbf{数で確かめる}\quad $x=2$ で、$\dfrac{\ln2.001-\ln2}{0.001}=0.49988\cdots$。たしかに $\dfrac{1}{2}$ に近い値です。

\bksubsection{8}{どんな指数でも $(x^r)'=rx^{r-1}$}

\label{bk:5-14:xr}5-6で、$(\sqrt{x})'$ を求めるのに、$(x^n)'=nx^{n-1}$ が分数の $n$ でも使えるとしました。その理由がここでわかります。$x>0$ で、$r$ をどんな数としても、⑥と同じように $x=e^{\ln x}$ なので $x^r=e^{r\ln x}$ です。内側を $u=r\ln x$（微分すると $r\times\dfrac{1}{x}$）とする合成関数の微分で、
\[ (x^r)'=e^{r\ln x}\times\Bigl(r\times\frac{1}{x}\Bigr)=x^r\times\frac{r}{x}=rx^{r-1} \]
最後は、$x^r\div x=x^{r-1}$（2-12の指数のきまり）を使いました。

\bksubsection{9}{$e^{ax}$ と $\frac{1}{x}$ の積分}

\label{bk:5-14:sekibun}微分の表を逆に読むと（5-11）、
\[ \int e^{ax}\,\dd x=\frac{1}{a}e^{ax}+C,\qquad \int\frac{1}{x}\,\dd x=\ln x+C\quad（x>0） \]
\noindent\textbf{なぜ $\dfrac{1}{a}$ がつくの？}\quad $e^{ax}$ を微分すると $a$ 倍になってしまうので、はじめに $\dfrac{1}{a}$ をかけておけば、微分したときにちょうどもとにもどります。微分して確かめます。
\[ \Bigl(\frac{1}{a}e^{ax}\Bigr)'=\frac{1}{a}\times a\,e^{ax}=e^{ax} \]
たとえば
\begin{align*}
a=2\ \text{なら}&\quad \int e^{2x}\,\dd x=\frac{1}{2}e^{2x}+C\\
a=-1\ \text{なら}&\quad \int e^{-x}\,\dd x=\frac{1}{-1}e^{-x}+C=-e^{-x}+C
\end{align*}
です。

\noindent\textbf{$\frac{1}{x}$ の積分}\par
5-11で「$n=-1$ のときは使えない」とした $\displaystyle\int x^{-1}\,\dd x$ の答えが、自然対数でした。⑦で $(\ln x)'=\dfrac{1}{x}$ だったからです。

\noindent\begin{minipage}[c]{0.44\linewidth}
\begin{tikzpicture}
\begin{axis}[bkplot, width=54mm, height=42mm, xmin=0, xmax=2.6, ymin=0, ymax=1.3, xtick={1,1.5,2}, ytick={0.5,1}, xlabel={$x$}, ylabel={$y$}]
  \foreach \k in {0,...,3} {\edef\tmp{\noexpand\fill[bkBlue!20, draw=bkBlue!60, line width=0.4pt] (axis cs:{1+0.25*\k},0) rectangle (axis cs:{1.25+0.25*\k},{1/(1.125+0.25*\k)});}\tmp}
  \addplot[bkcurve, domain=0.8:2.5] {1/x};
\end{axis}
\end{tikzpicture}
\end{minipage}\hfill
\begin{minipage}[c]{0.52\linewidth}\raggedright
面積で確かめます。$y=\dfrac{1}{x}$ の1から2までの下の面積を、はば0.25の長方形4つで近似します。高さは、区間のまん中（1.125、1.375、1.625、1.875）での $y$ の値です。
\end{minipage}

\noindent 4つの長方形の面積をたすと、
\[ \Bigl(\frac{1}{1.125}+\frac{1}{1.375}+\frac{1}{1.625}+\frac{1}{1.875}\Bigr)\times0.25=0.691\cdots \]
積分で求めると $\Bigl[\,\ln x\,\Bigr]_1^2=\ln2-\ln1=0.693\cdots-0=0.693\cdots$ で、ほとんど同じです。

\bksubsection{10}{「自分に比例してふえる」を表計算でつくる}

\label{bk:5-14:euler}$(e^x)'=e^x$ は、「ふえる速さが、いまの値と同じ」ということです。この性質だけを使って、5-8の方法（次の値＝今の値＋変化の速さ×$h$）で、$x=0$ から1まで計算してみましょう。はじめの値は1、変化の速さは「今の値」です。
\[ \text{次の値}=\text{今の値}+\text{今の値}\times h=\text{今の値}\times(1+h) \]
1回ごとに $(1+h)$ 倍になるので、$h=\dfrac{1}{n}$ で $n$ 回くり返すと、$x=1$ での値は $\Bigl(1+\dfrac{1}{n}\Bigr)^n$ です。

\begin{bkcenter}
\begin{bktable}{colspec={Q[c,m,wd=20mm] Q[c,m,wd=20mm] Q[c,m,wd=40mm]}, rowsep=1.5pt}
$h$ & 回数 $n$ & $x=1$ での値 $(1+h)^n$ \\
0.1 & 10 & 2.5937… \\
0.01 & 100 & 2.7048… \\
0.001 & 1000 & 2.7169… \\
\end{bktable}
\bktabcaption{$h$ を小さくすると、$e=2.71828\cdots$ に近づく。}
\end{bkcenter}
\bkhatenaref の銀行の計算（②）と、まったく同じ式になりました。「1年を $n$ 回に分けて、そのたびに今のお金に比例してふやす」ことは、「ふえる速さが自分に比例する」を少しずつ計算することだったのです。$n$ を大きくした極限で、ぴったり $e$ になります。

\noindent\begin{minipage}[c]{0.46\linewidth}
\begin{tikzpicture}
\begin{axis}[bkplot, width=54mm, height=46mm, xmin=0, xmax=1.2, ymin=0.8, ymax=3, xtick={0.25,0.5,0.75,1}, xticklabels={,0.5,,1}, ytick={1,2.44,2.72}, yticklabels={1,2.44,$e$}, xlabel={$x$}, ylabel={値}]
  \addplot[bkcurve, domain=0:1.1] {exp(x)};
  \addplot[bkcurve2, line width=1pt, mark=*, mark size=1.4pt] coordinates {(0,1) (0.25,1.25) (0.5,1.5625) (0.75,1.953125) (1,2.44140625)};
\end{axis}
\end{tikzpicture}
\end{minipage}\hfill
\begin{minipage}[c]{0.5\linewidth}\raggedright
オレンジは $h=0.25$（4回）で、1回ごとに1.25倍した折れ線です。$x=1$ で2.44になり、②の「3か月ごと」と同じ数です。青い $y=e^x$ より少し下にあります。$h$ を小さくすると、折れ線は青い曲線に近づきます。
\end{minipage}

\Needspace{14\baselineskip}
\bksubsection{11}{へっていく $e^{-t/\tau}$ と時定数}

\label{bk:5-14:tau}\noindent\begin{minipage}[c]{0.52\linewidth}
\begin{tikzpicture}
\begin{axis}[bkplot, width=64mm, height=46mm, xmin=0, xmax=5.4, ymin=0, ymax=1.15, xtick={1,2,3,4,5}, xticklabels={$\tau$,$2\tau$,$3\tau$,$4\tau$,$5\tau$}, ytick={0.368,1}, yticklabels={0.37,1}, xlabel={$t$}, ylabel={$\frac{N}{N_0}$}]
  \addplot[bkcurve, domain=0:5.2] {exp(-x)};
  \addplot[bkcurve2, line width=1pt, domain=0:1] {1-x};
  \draw[bkGray, dashed, line width=0.5pt] (axis cs:1,0) -- (axis cs:1,0.368) -- (axis cs:0,0.368);
\end{axis}
\end{tikzpicture}
\end{minipage}\hfill
\begin{minipage}[c]{0.44\linewidth}
へっていく量は、$N=N_0e^{-t/\tau}$ と書くことがよくあります。$\tau$（タウ）は時間の単位をもつ数で、\textbf{時定数}といいます。\bkidx{じていすう}{時定数 $\tau$}
\begin{itemize}[leftmargin=1.2\zw, itemsep=1pt, topsep=2pt, label={\color{bkTheme}●}]
\item $t=\tau$ で、$e^{-1}=0.368$ 倍（約37\,\%）になる。
\item $t=0$ での接線（オレンジ）は、ちょうど $t=\tau$ で0になる。
\item $5\tau$ たつと $e^{-5}=0.0067$ 倍で、ほぼ0。
\end{itemize}
\end{minipage}\par
\noindent\textbf{なぜ接線が $\tau$ で0になるの？}\quad ⑤で $a=-\dfrac{1}{\tau}$ とすると、$N'=N_0\times\Bigl(-\dfrac{1}{\tau}\Bigr)e^{-t/\tau}$ です。$t=0$ では $e^0=1$ なので、へる速さは $\dfrac{N_0}{\tau}$。はじめの量 $N_0$ を、この速さのままへらし続けると、$N_0\div\dfrac{N_0}{\tau}=\tau$ の時間で0になります。

「はじめの速さのままへり続けたら、$\tau$ で0になる」というのが、時定数の意味です。じっさいには、へるにつれてへる速さもおそくなるので、$\tau$ たっても37\,\%残ります。電気回路や計器の反応の速さは、時定数で表します（6-4）。

\begin{bkfigure}
\begin{tikzpicture}
\begin{axis}[bkplot, name=p1, at={(0mm,0mm)}, width=38mm, height=25mm, xmin=0, xmax=5.5, ymin=0, ymax=1.15, xtick=\empty, ytick=\empty, title={\sffamily\scriptsize 放射性物質の量}, title style={yshift=-1.5mm}]
  \addplot[bkcurve, line width=1.2pt, domain=0:5.3] {exp(-0.6*x)};
\end{axis}
\begin{axis}[bkplot, at={(45mm,0mm)}, name=p2, width=38mm, height=25mm, xmin=0, xmax=5.5, ymin=0, ymax=1.15, xtick=\empty, ytick=\empty, title={\sffamily\scriptsize お湯の温度}, title style={yshift=-1.5mm}]
  \addplot[bkcurve2, line width=1.2pt, domain=0:5.3] {0.25+0.75*exp(-0.7*x)};
  \addplot[bkthin, dashed, domain=0:5.3] {0.25};
\end{axis}
\begin{axis}[bkplot, at={(90mm,0mm)}, width=38mm, height=25mm, xmin=0, xmax=5.5, ymin=0, ymax=1.15, xtick=\empty, ytick=\empty, title={\sffamily\scriptsize コンデンサーの電圧}, title style={yshift=-1.5mm}]
  \addplot[bkcurve3, line width=1.2pt, domain=0:5.3] {1-exp(-0.9*x)};
  \addplot[bkthin, dashed, domain=0:5.3] {1};
\end{axis}
\end{tikzpicture}
\bkfigcaption{身のまわりの $e^{-t/\tau}$。横軸は時間。へる・冷める（差がへる）・たまる（足りない分がへる）が、同じ数学で書ける。}
\end{bkfigure}

\begin{column}{ネイピアと対数表}
$e$ の名前のもとになったスコットランドのネイピアは、400年ほど前、かけ算をたし算に変えて計算を楽にする「対数」を考え出し、20年かけて対数の表を作りました。天文学者のケプラーは、この表のおかげで惑星の計算を何年分も短くできたといわれます。$e$ そのものを使い始めたのは、その100年ほど後のスイスのオイラーで、$e$ という文字もオイラーが使い始めました。
\end{column}

\begin{senmonsho}[大学の教科書ふうの書き方]
\begin{bkhonbun}
指数関数 $\exp x\equiv e^x$ は、すべての実数 $x$ に対して $\dfrac{\dd}{\dd x}\exp x=\exp x$、$\exp0=1$ をみたす。その逆関数を $\ln x$ と書く。
\end{bkhonbun}
\bkyaku{$\exp x$（エクスポネンシャル・エックス）は $e^x$ の別の書き方。肩の式が長いときに使う：$e^{-t/\tau}=\exp(-t/\tau)$。「逆関数」は、入力と出力を入れかえた関数（$y=e^x$ と $x=e^y$）。}
\end{senmonsho}

\bkkigocards{\bkkigocard{$e$}{イー（ネイピア数）}{$\bigl(1+\frac{1}{n}\bigr)^n$ の近づく先。$2.718\cdots$。$(e^x)'=e^x$}{Napier's constant}{5-14}}
{\bkkigocard{$\ln x$}{エル・エヌ・エックス}{自然対数。$e$ を何乗すると $x$ になるか}{natural log}{5-14}}

\begin{taisetsu}
「1年あたり100\,\%」を $n$ 回に分けると、1回 $\dfrac{1}{n}$ ずつ、1年で $\Bigl(1+\dfrac{1}{n}\Bigr)^n$ 倍 → $e$ に近づく\\
$(e^x)'=e^x$\qquad $(e^{ax})'=a\,e^{ax}$\qquad $(\ln x)'=\dfrac{1}{x}$\qquad $(a^x)'=a^x\ln a$\\
$\displaystyle\int e^{ax}\,\dd x=\frac{1}{a}e^{ax}+C$\qquad $\displaystyle\int\frac{1}{x}\,\dd x=\ln x+C$
\end{taisetsu}

\begin{kotoba}
\small
\bkword{複利}{利息にも利息がつく、ふえ方。}
\bkword{$e$（ネイピア数）}{$(1+\frac{1}{n})^n$ の近づく先。2.71828…。$e^x$ は微分しても変わらない。}
\bkword{自然対数 $\ln x$}{$e$ を底とする対数。$\log_e x$。}
\bkword{$\exp x$}{$e^x$ の別の書き方。}
\end{kotoba}

%% ---------------- 例題 ----------------
\bksection{reidai}{例題でたしかめよう}

\begin{reidai}[利息を何回かに分ける]
「1年あたり100\,\%」の利率で1万円を預け、利息を3か月ごと（1年に4回）に計算してもらいます。1年後のお金を求めましょう。\par
\kangaekata 3か月は1年の $\dfrac{1}{4}$ なので、1回分の利率は $100\,\%\times\dfrac{1}{4}=25\,\%$ です（①②）。毎回、そのときのお金の25\,\%がつくので、1回ごとに1.25倍になります。\par
\shiki 1回分の利率\quad $100\,\%\times\dfrac{1}{4}=25\,\%$、1回ごとに $1+0.25=1.25$ 倍\par
\hspace*{2.6\zw}4回くり返すので\quad $1\times1.25^4=1.25\times1.25\times1.25\times1.25$\par
\hspace*{2.6\zw}$=2.44140625$\par
\kotae 約2.44万円（1年に1回の2万円、半年ごとの2.25万円より多い）
\end{reidai}

\begin{jibun}
「1年あたり100\,\%」の利率で1万円を預け、利息を4か月ごと（1年に3回）に計算してもらいます。1年後のお金を求めましょう。分数のまま計算してから、小数にしましょう。
\bkjans{1回分 $\dfrac{1}{3}$、$\Bigl(1+\dfrac{1}{3}\Bigr)^3=\Bigl(\dfrac{4}{3}\Bigr)^3=\dfrac{64}{27}\approx2.37$。約2.37万円}
\end{jibun}

\begin{reidai}[$e$ をふくむ関数の微分]
次の関数を微分しましょう。\quad（1）$e^{4x}$\quad（2）$3e^{-x}$\quad（3）$e^{x^2}$\par
\kangaekata（1）（2）は⑤の $(e^{ax})'=ae^{ax}$（（2）は $a=-1$）。（3）は内側 $u=x^2$ の合成関数（5-6）。\par
\shiki（1）$4e^{4x}$\quad（2）$3\times(-1)e^{-x}=-3e^{-x}$\par
\hspace*{2.6\zw}（3）外側 $e^u$ の微分 $e^u$ に、内側 $x^2$ の微分 $2x$ をかけて $e^{x^2}\times2x$\par
\kotae （1）$4e^{4x}$\quad（2）$-3e^{-x}$\quad（3）$2xe^{x^2}$
\end{reidai}

\begin{jibun}
次の関数を微分しましょう。\quad（1）$e^{5x}$\quad（2）$2e^{-3x}$\quad（3）$e^{2x+1}$
\bkjans{（1）\nobreak $5e^{5x}$\quad（2）\nobreak $-6e^{-3x}$\quad（3）\nobreak $2e^{2x+1}$}
\end{jibun}

\begin{reidai}[へる速さ]
ある物質の量が $N=1000e^{-0.1t}$（g、$t$ は時間）でへっていきます。はじめ（$t=0$）のへる速さを求めましょう。\par
\kangaekata へる速さは、$N$ を $t$ で微分した $\dfrac{\dd N}{\dd t}$。⑤の $(e^{ax})'=ae^{ax}$ で、$a=-0.1$ です。\par
\shiki $\dfrac{\dd N}{\dd t}=1000\times(-0.1)e^{-0.1t}=-100e^{-0.1t}$\par
\hspace*{2.6\zw}$t=0$ では $e^0=1$ なので $-100\times1=-100$\par
\kotae 1時間あたり\qty{100}{g}の速さでへる（$\frac{\dd N}{\dd t}=-0.1N$：へる速さは、いまの量の0.1倍）
\end{reidai}

\begin{jibun}
$N=800e^{-0.05t}$ について、$\dfrac{\dd N}{\dd t}$ を求め、$t=0$ での値を答えましょう。
\bkjans{$800\times(-0.05)e^{-0.05t}=-40e^{-0.05t}$、$t=0$ で $-40$}
\end{jibun}

\begin{reidai}[$\ln$ と $a^x$ の微分]
次の関数を微分しましょう。\quad（1）$\ln(2x+1)$\quad（2）$2^x$\par
\kangaekata（1）内側を $u=2x+1$ とする合成関数。外側 $\ln u$ の微分は $\dfrac{1}{u}$（⑦）。（2）$2^x=e^{x\ln2}$ に直す（⑥）。\par
\shiki（1）$\dfrac{1}{2x+1}\times(2x+1)'=\dfrac{1}{2x+1}\times2$\par
\hspace*{2.6\zw}（2）$(2^x)'=(e^{x\ln2})'=e^{x\ln2}\times\ln2=2^x\ln2$\par
\kotae （1）$\dfrac{2}{2x+1}$\quad（2）$2^x\ln2$（$\approx0.693\times2^x$）
\end{reidai}

\begin{jibun}
次の関数を微分しましょう。\quad（1）$\ln(x^2+1)$\quad（2）$10^x$
\bkjans{（1）\nobreak $\dfrac{2x}{x^2+1}$\quad（2）\nobreak $10^x\ln10$}
\end{jibun}

\begin{reidai}[$e^{ax}$ と $\frac{1}{x}$ の積分]
（1）$\displaystyle\int_0^1e^x\,\dd x$\quad（2）$\displaystyle\int_0^1e^{-x}\,\dd x$\quad（3）$\displaystyle\int_1^e\frac{1}{x}\,\dd x$ を求めましょう。\par
\kangaekata ⑨の積分の式で原始関数を求め、上端の値から下端の値をひきます（5-11）。（2）は $a=-1$ なので、$\dfrac{1}{a}=-1$ がつきます。\par
\shiki（1）$\Bigl[\,e^x\,\Bigr]_0^1=e^1-e^0=e-1$\par
\hspace*{2.6\zw}（2）$\Bigl[\,-e^{-x}\,\Bigr]_0^1=(-e^{-1})-(-e^0)=1-e^{-1}$\par
\hspace*{2.6\zw}（3）$\Bigl[\,\ln x\,\Bigr]_1^e=\ln e-\ln1=1-0$\par
\kotae （1）$e-1\approx1.718$\quad（2）$1-e^{-1}\approx0.632$\quad（3）1
\end{reidai}

\begin{jibun}
（1）$\displaystyle\int_0^2e^{x}\,\dd x$\quad（2）$\displaystyle\int_1^3\frac{1}{x}\,\dd x$\quad（3）$\displaystyle\int_0^1e^{2x}\,\dd x$ を求めましょう。
\bkjans{（1）\nobreak $e^2-1\approx6.39$\quad（2）\nobreak $\ln3\approx1.10$\quad（3）\nobreak $\dfrac{e^2-1}{2}\approx3.19$}
\end{jibun}

\begin{reidai}[接線で近似する]
$y=e^x$ のグラフの、$x=0$ の点での接線の式を求めましょう。これを使って $e^{0.2}$ の近似値を求め、電卓の値 1.2214… とくらべましょう。\par
\kangaekata 接線の式は $y=f(a)+f'(a)(x-a)$（5-8）。$f(x)=e^x$、$a=0$ です。\par
\shiki $f(0)=e^0=1$、$f'(x)=e^x$ より $f'(0)=1$\par
\hspace*{2.6\zw}接線は $y=1+1\times(x-0)=x+1$\par
\hspace*{2.6\zw}$x=0.2$ を入れて $e^{0.2}\approx0.2+1=1.2$\par
\kotae 接線 $y=x+1$、$e^{0.2}\approx1.2$（ずれは約0.02）
\end{reidai}

\begin{jibun}
例題6の接線を使って、$e^{-0.1}$ の近似値を求め、電卓の値 0.9048… とくらべましょう。
\bkjans{$e^{-0.1}\approx-0.1+1=0.9$（ずれは約0.005）}
\end{jibun}

\begin{reidai}[半減期]
$N=N_0e^{-0.1t}$（$t$ は時間）でへる物質が、はじめの半分になる時間（半減期）を求めましょう。\par\bkidx{はんげんき}{半減期}
\kangaekata 半分になるのは、$N_0e^{-0.1T}=\dfrac{1}{2}N_0$、つまり $e^{-0.1T}=\dfrac{1}{2}$ のとき。$T$ が指数にあるので、両辺の自然対数をとって $T$ を下ろします（⑥の $\ln e^y=y$）。\par
\shiki 両辺の自然対数をとって\quad $\ln e^{-0.1T}=\ln\dfrac{1}{2}$\par
\hspace*{2.6\zw}左辺は $-0.1T$。右辺は $\ln\dfrac{1}{2}=\ln2^{-1}=-\ln2$（3-12）\par
\hspace*{2.6\zw}だから $-0.1T=-\ln2$\par
\hspace*{2.6\zw}$T=\dfrac{\ln2}{0.1}=\dfrac{0.693}{0.1}=6.93$\par
\kotae 約6.9時間
\end{reidai}

\begin{jibun}
$N=N_0e^{-0.05t}$（$t$ は年）でへる物質の半減期を求めましょう。
\bkjans{\mbox{$\dfrac{\ln2}{0.05}\approx13.9$\,年}}
\end{jibun}

\begin{reidai}[時定数]
コンデンサーにたまった電気の量が、$Q=Q_0e^{-t/2}$（$t$ は秒）でへっていきます。時定数と、半分になるまでの時間を求めましょう。\par
\kangaekata $e^{-t/\tau}$ とくらべて $\tau$ を読む（⑪）。半分の時間は、例題7と同じく両辺の自然対数をとる。\par
\shiki $e^{-t/2}$ なので時定数 $\tau=2$ 秒\par
\hspace*{2.6\zw}半分になる時間 $T$ は $e^{-T/2}=\dfrac{1}{2}$\par
\hspace*{2.6\zw}$-\dfrac{T}{2}=\ln\dfrac{1}{2}=-\ln2$ より $T=2\ln2\approx2\times0.693\approx1.39$\par
\kotae 時定数2秒、半分になるのは約1.4秒後（半減期＝時定数×$\ln2$）
\end{reidai}

\begin{jibun}
温度計の針が、$\theta=\theta_0e^{-t/5}$（$t$ は秒）で目もりに近づいていきます。時定数と、ずれが半分になるまでの時間を求めましょう。
\bkjans{時定数5秒、\mbox{$5\ln2\approx3.5$\,秒}}
\end{jibun}

\begin{reidai}[たえずふえる量]
いまある量に比例して、たえずふえていく量は $y=y_0e^{rt}$ と書けます（$r$ はふえる割合。$(y_0e^{rt})'=ry_0e^{rt}=ry$ なので、ふえる速さはいまの量の $r$ 倍）。100万円が $r=0.02$（1年あたり）でたえずふえるとき、10年後の金額を求めましょう。\par
\kangaekata $y=y_0e^{rt}$ に、$y_0=100$（万円）、$r=0.02$、$t=10$ を入れます。\par
\shiki $100\times e^{0.02\times10}=100\times e^{0.2}=100\times1.2214=122.14$\par
\kotae 約122万円（電卓で $e^{0.2}$ を計算する。表計算では \texttt{EXP(0.2)}）
\end{reidai}

\begin{jibun}
人口が $r=0.05$（1年あたり）でたえずふえる町があります。いま1万人のとき、10年後の人口を求めましょう（$e^{0.5}=1.6487$）。
\bkjans{$10000\times e^{0.5}=16487$。約16487人}
\end{jibun}

\begin{reidai}[コンデンサーにたまる電圧]
電池につないだコンデンサーの電圧が、$V=10\,(1-e^{-t/2})$（V、$t$ は秒）でふえていきます。2秒後の電圧と、$t=0$ での電圧のふえる速さを求めましょう（$e^{-1}=0.368$）。\par
\kangaekata $1-e^{-t/2}$ は、1から「へっていく量」をひいた形。時定数 $\tau=2$ 秒たつと、最後の値の $1-0.368=0.632$（約63\,\%）までたまります。\par
\shiki $t=2$：$10\times(1-e^{-1})=10\times(1-0.368)=6.32$\par
\hspace*{2.6\zw}$\dfrac{\dd V}{\dd t}=10\times\Bigl(-\Bigl(-\dfrac{1}{2}\Bigr)e^{-t/2}\Bigr)=5e^{-t/2}$、$t=0$ で5\par
\kotae 約\qty{6.3}{V}、毎秒\qty{5}{V}（だんだんゆっくりになり、\qty{10}{V}に近づく）
\end{reidai}

\begin{jibun}
コンデンサーの電圧が $V=5\,(1-e^{-t/3})$（V）でふえていきます。3秒後の電圧と、$t=0$ でのふえる速さを求めましょう（$e^{-1}=0.368$）。
\bkjans{$5\times0.632\approx\qty{3.16}{V}$。$\dfrac{\dd V}{\dd t}=\dfrac{5}{3}e^{-t/3}$、$t=0$ で毎秒約\qty{1.67}{V}}
\end{jibun}

\begin{reidai}[2倍になるまでの時間]
細菌の数が $N=50e^{0.35t}$（$t$ は時間）でふえます。数が2倍になるまでの時間（倍加時間）と、$t=0$ でのふえる速さを求めましょう（$\ln2=0.693$）。\par
\kangaekata 半減期（例題7）と同じように、$e^{0.35T}=2$ の両辺の自然対数をとります。\par
\shiki $0.35T=\ln2$ より $T=\dfrac{0.693}{0.35}\approx1.98$\par
\hspace*{2.6\zw}$\dfrac{\dd N}{\dd t}=50\times0.35e^{0.35t}=17.5e^{0.35t}$、$t=0$ で17.5\par
\kotae 約2時間。1時間あたり17.5（いまの数の0.35倍の速さでふえる）
\end{reidai}

\begin{jibun}
$N=200e^{0.1t}$（$t$ は時間）でふえる量の、倍加時間と、$t=0$ でのふえる速さを求めましょう。
\bkjans{\mbox{$T=\dfrac{0.693}{0.1}\approx6.9$\,時間}。$\dfrac{\dd N}{\dd t}=20e^{0.1t}$、$t=0$ で1時間あたり20}
\end{jibun}

\begin{reidai}[$\frac{1}{x}$ の積分と、対数のきまり]
$\displaystyle\int_2^6\frac{1}{x}\,\dd x$ を求め、$\ln$ 1つで表しましょう。\par
\kangaekata ⑨で原始関数は $\ln x$。上端の値から下端の値をひくと $\ln6-\ln2$ になります。対数のきまり $\ln(bc)=\ln b+\ln c$（3-12）を使うと、$\ln6=\ln(3\times2)=\ln3+\ln2$ なので、ひき算がかんたんになります。\par
\shiki $\displaystyle\int_2^6\frac{1}{x}\,\dd x=\Bigl[\,\ln x\,\Bigr]_2^6=\ln6-\ln2$\par
\hspace*{2.6\zw}$\ln6=\ln(3\times2)=\ln3+\ln2$ なので\par
\hspace*{2.6\zw}$\ln6-\ln2=(\ln3+\ln2)-\ln2=\ln3$\par
\kotae $\ln3$（$\approx1.10$。$6\div2=3$ の $\ln$ になる）
\end{reidai}

\begin{jibun}
$\displaystyle\int_5^{10}\frac{1}{x}\,\dd x$ を求め、$\ln$ 1つで表しましょう。
\bkjans{$\ln10-\ln5=(\ln2+\ln5)-\ln5=\ln2\approx0.693$}
\end{jibun}

\begin{reidai}[ふつうの利率を、細かく分ける]
1年あたり6\,\%の利率で、100万円を預けます。（1）1年に1回、（2）毎月（1年に12回）利息を計算するとき、1年後の金額を求めましょう。（3）たえず（いつも少しずつ）ふえるときの金額 $100e^{0.06}$ とくらべましょう（$e^{0.06}=1.06184$）。\par
\kangaekata 1回分の利率は、①②と同じく「1年あたりの利率×1回分の時間（年）」です。毎月なら1か月は $\dfrac{1}{12}$ 年なので、$6\,\%\times\dfrac{1}{12}=0.5\,\%$。1回ごとに $1+0.005=1.005$ 倍になり、それを12回くり返します。\par
\shiki（1）$100\times1.06=106$\par
\hspace*{2.6\zw}（2）$100\times1.005^{12}=100\times1.06168=106.168$\par
\hspace*{2.6\zw}（3）$100\times1.06184=106.184$\par
\kotae （1）106万円\quad（2）約106.17万円\quad（3）約106.18万円（細かく分けるほど、$y_0e^{rt}$ の値に近づく）
\end{reidai}

\noindent 1年あたりの利率を $r$ として $n$ 回に分けると、1年後は $\Bigl(1+\dfrac{r}{n}\Bigr)^n$ 倍です。$n$ を大きくすると、これは $e^r$ に近づきます。\bkreidairef{9}の「たえずふえる量」$y_0e^{rt}$ は、この行き着く先でした（$r=1$ のときが②の $e$）。

\begin{jibun}
1年あたり12\,\%の利率で100万円を預けます。毎月（1年に12回）利息を計算するとき、1回分の利率と、1年後の金額を求めましょう（$1.01^{12}=1.12683$）。
\bkjans{1回分 $12\,\%\times\dfrac{1}{12}=1\,\%$。$100\times1.01^{12}\approx112.68$。約112.68万円}
\end{jibun}

%% ---------------- まちがえやすいところ・物理工学 ----------------
\bksection{tsumazuki}{まちがえやすいところをチェック}

\begin{tsumazuki}{半年ごとの利率を、2回めだけ100\,\%にしてしまう}
\bkng{はじめの半年は50\,\%、次の半年は100\,\%で、$1.5\times2=3$ 万円}\bkok{半年は毎回50\,\%で、$1.5\times1.5=2.25$ 万円}
100\,\%は「1年あたり」の割合です。半年はいつでも1年の半分なので、何回めでも50\,\%です（①）。
\end{tsumazuki}

\begin{tsumazuki}{$e^x$ に $(x^n)'=nx^{n-1}$ を使ってしまう}
\bkng{$(e^x)'=xe^{x-1}$}\bkok{$(e^x)'=e^x$}
$x^n$ は「$x$ を $n$ 回かける」で、変わるのは底の $x$。$e^x$ は「$e$ を $x$ 回かける」で、変わるのは指数の $x$。形は似ていても、まったくちがう関数です。
\end{tsumazuki}

\begin{tsumazuki}{$\ln(a+b)=\ln a+\ln b$ としてしまう}
対数で和になるのは、かけ算のときです：$\ln(ab)=\ln a+\ln b$（3-12）。$\ln(a+b)$ は、それ以上かんたんにはなりません。$\ln(1+1)=\ln2=0.693$ で、$\ln1+\ln1=0$ ではありません。
\end{tsumazuki}

\begin{tsumazuki}{時定数と半減期をまちがえる}
時定数 $\tau$ は「約37\,\%（$e^{-1}$）になる時間」、半減期は「50\,\%になる時間（$\tau\ln2\approx0.69\tau$）」で、ちがう量です。
\end{tsumazuki}

\begin{butsuri}
\begin{itemize}[leftmargin=1.2\zw, itemsep=1pt, topsep=0pt, label={\color{bkPurple}●}]
\item \textbf{放射線と年代測定}\\炭素14は、半減期約5730年で、$e^{-\lambda t}$ の形でへります。木や骨に残っている量から、何年前のものかを調べられます。
\item \textbf{冷めるお湯・コンデンサー}\\お湯と部屋の温度差や、コンデンサーにたまった電気は、$e^{-t/\tau}$ の形でへります。時定数 $\tau$ が、計器や制御の速さを表します（6-4）。
\item \textbf{ラプラス変換}\\制御工学で使うラプラス変換は、$e^{-st}$ をかけて積分する計算です。$e$ の微分・積分が、その土台になります（9-5）。
\end{itemize}
\end{butsuri}


%% ---------------- 章末問題 ----------------
\bkendproblems
\bkproblemgroup{基本}{3}

\bkproblem 「1年あたり100\,\%」の利率で1万円を預け、利息を2か月ごと（1年に6回）に計算してもらいます。1回分の利率は何\,\%ですか。1年後のお金は約何万円ですか（小数第2位まで）。
\bkpans{$\dfrac{100}{6}\approx16.7$\,\%。$\Bigl(1+\dfrac{1}{6}\Bigr)^6\approx2.52$ 万円}

\bkproblem 次の関数を微分しましょう。\quad（1）$e^{6x}$\quad（2）$4e^{-2x}$\quad（3）$e^{x^3}$
\bkpans{（1）\nobreak $6e^{6x}$\quad（2）\nobreak $-8e^{-2x}$\quad（3）\nobreak $3x^2e^{x^3}$}

\bkproblem 次の関数を微分しましょう。\quad（1）$\ln(3x)$\quad（2）$\ln(4x-1)$\quad（3）$3^x$
\bkpans{（1）\nobreak $\dfrac{3}{3x}=\dfrac{1}{x}$\quad（2）\nobreak $\dfrac{4}{4x-1}$\quad（3）\nobreak $3^x\ln3$}

\bkproblem 次の定積分を求めましょう。\quad（1）$\displaystyle\int_0^2e^{-x}\,\dd x$\quad（2）$\displaystyle\int_2^4\frac{1}{x}\,\dd x$
\bkpans{（1）\nobreak $1-e^{-2}\approx0.865$\quad（2）\nobreak $\ln4-\ln2=\ln2\approx0.693$}

\bkproblem $N=500e^{-0.2t}$ について、$\dfrac{\dd N}{\dd t}$ と半減期を求めましょう。
\bkpans{$-100e^{-0.2t}$。半減期 $\dfrac{\ln2}{0.2}\approx3.5$}

\bkproblem コンデンサーの電圧が $V=12\,(1-e^{-t/4})$（V）でふえていきます。4秒後の電圧と、$t=0$ での電圧のふえる速さを求めましょう（$e^{-1}=0.368$）。
\bkpans{$12\times0.632\approx\qty{7.6}{V}$。$\dfrac{\dd V}{\dd t}=3e^{-t/4}$、$t=0$ で毎秒\qty{3}{V}}

\bkproblem 電気の量が $Q=Q_0e^{-t/4}$（$t$ は秒）でへっていきます。時定数と、半分になるまでの時間を求めましょう。
\bkpans{時定数4秒、\mbox{$4\ln2\approx2.8$\,秒}}

\bkproblem 「1年あたり100\,\%」の利率で1万円を預け、利息を毎月（1年に12回）計算してもらいます。1回分の利率は1年の何分の1ですか。1年後のお金は約何万円ですか（電卓で小数第2位まで）。
\bkpans{$\dfrac{1}{12}$。$\Bigl(1+\dfrac{1}{12}\Bigr)^{12}\approx2.61$ 万円}

\bkproblem 次の関数を微分しましょう。\quad（1）$5e^{2x}$\quad（2）$e^{-x^2}$\quad（3）$e^{3x+2}$
\bkpans{（1）\nobreak $10e^{2x}$\quad（2）\nobreak $-2xe^{-x^2}$\quad（3）\nobreak $3e^{3x+2}$}

\bkproblem 100万円が $r=0.03$（1年あたり）でたえずふえるとき、10年後の金額を求めましょう（$e^{0.3}=1.3499$）。
\bkpans{$100\times e^{0.3}=134.99$。約135万円}

\bkproblemgroup{チャレンジ}{4}

\bkproblem $y=e^x$ のグラフの、$x=0$ の点での接線の式を求めましょう。これを使って、$e^{0.1}$ の近似値を求め、電卓の値 1.10517… とくらべましょう。
\bkpans{$y=x+1$。$e^{0.1}\approx1.1$（ずれ約0.005）}

\bkproblem 細菌の数が、$N=100e^{0.7t}$（$t$ は時間）でふえていきます。数が2倍になるまでの時間（倍加時間）と、$t=0$ でのふえる速さを求めましょう。
\bkpans{$e^{0.7T}=2$ より \mbox{$T=\dfrac{\ln2}{0.7}\approx1.0$\,時間}。$\dfrac{\dd N}{\dd t}=70e^{0.7t}$、$t=0$ で1時間あたり70}

\bkproblem 人口が $r=0.07$（1年あたり）でたえずふえる町があります。人口が2倍になるまでの年数を求めましょう（$\ln2=0.693$）。
\bkpans{$e^{0.07T}=2$ より \mbox{$T=\dfrac{0.693}{0.07}=9.9$\,年}}

\bkproblem $\displaystyle\int_0^{\ln2}e^x\,\dd x$ を求めましょう（$e^{\ln2}=2$ を使います）。
\bkpans{$\Bigl[\,e^x\,\Bigr]_0^{\ln2}=e^{\ln2}-e^0=2-1=1$}

\begin{furikaeri}
\begin{checklist}
\item $e$ が $(1+\frac{1}{n})^n$ の近づく先であることを確かめられる。\hfill{\footnotesize\sffamily\color{bkTheme}→\pageref{bk:5-14:n}ページ}
\item $(e^{ax})'=ae^{ax}$ を使える。\hfill{\footnotesize\sffamily\color{bkTheme}→\pageref{bk:5-14:eax}ページ}
\item $\ln x$ を微分・積分できる。\hfill{\footnotesize\sffamily\color{bkTheme}→\pageref{bk:5-14:lnbibun}ページ}
\end{checklist}
\end{furikaeri}

\bkfukidashi{sora}{ふえる・へるは $e$ で書けた。じゃあ、ゆれる・回るは？　次は $\sin$ と $\cos$ だ！}
